% 第九卷 7.6.3.3 节：语言状态上蒙特卡洛树搜索的一次迭代（跨栏，约160 mm宽）。
% 四个面板依次为选择、扩展、评估、回传；节点为“题目加已生成步骤”，第二行为(N, Q)。
% 数值为教学示例：根(3, 0.60)，子节点“书折后60元”(2, 0.80)、“全打八折”(1, 0.20)，
% 根的Q为两子节点按访问次数加权的平均；新步骤“笔20元”评估得0.9，
% 回传后“书折后60元”变为(3, 0.83)，根变为(4, 0.68)。
% 蓝色：状态节点；红色：本次迭代经过或更新的节点与路径；灰色虚线：新扩展而未评估的节点。
\begin{tikzpicture}[
  x=1mm,y=1mm,
  figure picture,
  st/.style={figure node,draw=FigureInput,fill=FigureInputFill,minimum height=9mm,
    minimum width=15mm,inner xsep=1.5pt,inner ysep=1.5pt},
  fresh/.style={st,draw=FigureMuted,fill=white,dash pattern=on 2pt off 1.5pt},
  hot/.style={st,draw=FigureOutput,fill=FigureOutputFill},
  path/.style={figure flow,draw=FigureOutput,line width=1pt},
  edge/.style={figure link}
]
  \foreach \i/\t in {0/选择,1/扩展,2/评估,3/回传}{
    \draw[figure subgroup] ({\i*41},-7) rectangle ({\i*41+38},46);
    \node[figure heading,anchor=north west] at ({\i*41+1.5},45) {\t};
  }

  % ---------- 面板一：选择 ----------
  \begin{scope}[shift={(0,0)}]
    \node[hot] (r) at (19,32) {题目\\(3, 0.60)};
    \node[hot] (a) at (10,17) {书折后60元\\(2, 0.80)};
    \node[st] (b) at (29,17) {全打八折\\(1, 0.20)};
    \draw[path] (r) -- (a);
    \draw[edge] (r) -- (b);
    \node[figure note] at (26,42) {按PUCT下行};
  \end{scope}

  % ---------- 面板二：扩展 ----------
  \begin{scope}[shift={(41,0)}]
    \node[st] (r) at (19,32) {题目\\(3, 0.60)};
    \node[st] (a) at (10,17) {书折后60元\\(2, 0.80)};
    \node[st] (b) at (29,17) {全打八折\\(1, 0.20)};
    \node[fresh,minimum width=13mm] (c) at (7.5,2) {笔20元\\(0, --)};
    \node[fresh,minimum width=13mm] (d) at (22,2) {笔16元\\(0, --)};
    \draw[edge] (r) -- (a);
    \draw[edge] (r) -- (b);
    \draw[edge,dashed] (a) -- (c);
    \draw[edge,dashed] (a) -- (d);
    \node[figure note] at (27,42) {采样2个候选};
  \end{scope}

  % ---------- 面板三：评估 ----------
  \begin{scope}[shift={(82,0)}]
    \node[st] (r) at (19,32) {题目\\(3, 0.60)};
    \node[st] (a) at (10,17) {书折后60元\\(2, 0.80)};
    \node[st] (b) at (29,17) {全打八折\\(1, 0.20)};
    \node[hot,minimum width=13mm] (c) at (7.5,2) {笔20元\\$v=0.9$};
    \node[fresh,minimum width=13mm] (d) at (22,2) {笔16元\\(0, --)};
    \draw[edge] (r) -- (a);
    \draw[edge] (r) -- (b);
    \draw[edge] (a) -- (c);
    \draw[edge,dashed] (a) -- (d);
    \node[figure note] at (25.5,42) {PRM或续写打分};
  \end{scope}

  % ---------- 面板四：回传 ----------
  \begin{scope}[shift={(123,0)}]
    \node[hot] (r) at (19,32) {题目\\(4, 0.68)};
    \node[hot] (a) at (10,17) {书折后60元\\(3, 0.83)};
    \node[st] (b) at (29,17) {全打八折\\(1, 0.20)};
    \node[hot,minimum width=13mm] (c) at (7.5,2) {笔20元\\(1, 0.90)};
    \node[fresh,minimum width=13mm] (d) at (22,2) {笔16元\\(0, --)};
    \draw[path] (c) -- (a);
    \draw[path] (a) -- (r);
    \draw[edge] (r) -- (b);
    \draw[edge,dashed] (a) -- (d);
    \node[figure note] at (27,42) {更新$N$与$Q$};
  \end{scope}
\end{tikzpicture}
